\subsection{Original blowup psh potentials give genuine local directed limits}

\begin{lemma}[Dyadic nets on the actual original blowup]
\label{lem:fc-blowup-dyadic-finite-fs-net}
\lean{EnvelopeProperty.ProjSubvariety.BlowupCenter.IsPSH.exists_isFinFS_net_dyadic}
\leanok
\uses{thm:fc-finite-fs-dyadic-net}
Every original psh potential of rational bidegree $(a,b)$ on the original
analytic blowup is the decreasing infimum of a genuine directed finite
Fubini--Study net of bidegrees $(a+2^{-n(\alpha)},b+2^{-n(\alpha)})$, with an
increasing cofinal natural index $n$. No countability of the field is required.
\end{lemma}
\begin{proof}
Apply the genuine dyadic-net regularization theorem to the actual original
blowup Fubini--Study system and the rational class $(a,b)$. Its system and
classes are exactly those defining the original blowup psh potential.
\end{proof}

\begin{lemma}[Normalized original chart potentials are directed convex PL limits]
\label{lem:fc-normalized-chart-convex-pl-net}
\lean{EnvelopeProperty.ProjSubvariety.BlowupCenter.fc_exists_convexPL_net_chartPotential}
\leanok
\uses{lem:fc-blowup-dyadic-finite-fs-net}
For every original blowup psh potential of class $(a,0)$, the actual normalized
potential on each genuine nonzero Rees chart over an integral original
affine base chart is the decreasing infimum of a
genuine directed net of original regular-unit convex PL functions. No
countability or prescribed approximation data are assumptions.
\end{lemma}
\begin{proof}
Write each of the constructed finite Fubini--Study approximants on the actual
Rees chart by its original finite rational terms. The unchanged chart formula
adds the bidegree normalization to its value. The two logarithmic chart scales
are nonnegative. The approximants decrease, and their dyadic classes decrease,
so these actual finite chart maxima decrease as well. Add the actual finite
standard truncation at the cofinal natural index to make every chart term a
finite original regular-unit convex PL function. These truncations decrease
to minus infinity and hence do not change the directed infimum.

The first logarithmic coefficient tends to $a$ and the second tends to zero.
Continuity of the finite real coefficient expression identifies the infimum
of the normalization terms with the unchanged original normalization.
Addition commutes with these directed decreasing infima because each
approximant and normalization term is finite above. The exact original
finite-term formula therefore identifies the total infimum with the original
normalized chart potential, retaining every original point and value.
\end{proof}

\begin{remark}[Relation to the source local psh class]
\label{rmk:fc-directed-pl-source-relation}
The source definition of locally psh uses decreasing sequences. The directed
limit statement above is kept separate from that definition. Its curve
comparison consequences follow directly from finite multiplicity-weighted
norms commuting with directed decreasing infima and the directed-limit
closure of source line subharmonic comparison. It does not assert that every
directed net has a cofinal sequence.
\end{remark}
